\documentclass[10pt,a4paper]{amsart}
\usepackage[margin=19mm]{geometry}
\usepackage{amsmath,amssymb,mathtools}
\usepackage[scr=rsfs,cal=euler]{mathalfa}
\usepackage{xfrac}
\usepackage[unicode,hidelinks,bookmarks=false]{hyperref}
\newcommand{\ie}{i.e.\ }
\newcommand{\cf}{cf.\ }
\newcommand{\resp}{respectively\ }
\newcommand{\Subsection}{\subsection}
\providecommand{\nobreakdash}{\nobreak\hbox{-}}
\setlength{\emergencystretch}{2em}
\multlinegap0pt
\usepackage{bm}
\usepackage{bbm}
\usepackage{dsfont}
\usepackage{xspace}
\usepackage{cancel}
\usepackage{mathtools}
\usepackage{amsthm}
\makeatletter
\@ifundefined{theorem}{\newtheorem{theorem}{Theorem}}{}
\@ifundefined{lemma}{\newtheorem{lemma}[theorem]{Lemma}}{}
\makeatother
\title[Persistent Homology for High-dimensional Data Based on Spect]{Persistent Homology for High-dimensional Data Based on Spectral Methods}
\author{Sebastian Damrich, Philipp Berens, Dmitry Kobak}
\date{Source: arXiv:2311.03087v3 (CC BY 4.0)}
\begin{document}
\maketitle

\noindent\emph{Source and licence.} Persistent Homology for High-dimensional Data Based on Spectral Methods. Version arXiv:2311.03087v3 (CC BY 4.0), distributed under the Creative Commons Attribution 4.0 International licence, as recorded in the arXiv licence metadata for this exact version and shown on the abstract page https://arxiv.org/abs/2311.03087v3. Full rights notice: licenses/arxiv-2311.03087-CC-BY-4.0.txt.\par\smallskip

\noindent\textit{Editorial scope.} Counted unit: the Lemma whose source label is \texttt{lem:pm} in the article (source0.tex lines 570--572, source label \texttt{lem:pm}), delivered together with the article's own complete original proof of it (source0.tex lines 573--616, one \texttt{proof} environment of 44 lines). No mathematics is written, repaired or re-derived: statement and proof are the article's own text line for line. Required context retained uncounted: none. Every definition, display and label of those lines is retained unchanged. Omitted from this package and not redistributed: 1--569, 617--1523. Nothing inside a retained excerpt is omitted or altered: every retained body is delivered exactly as the article's lines are written, so no omission receipt of any kind accompanies this package. Citations: the retained excerpts carry no citation command at all, so no bibliography entry of the article is transcribed and no reference is redistributed. Reference adaptation: none was needed and none was made; every reference inside the delivered unit resolves inside it, so no unresolved reference or citation is produced. Printed slips kept literal: no printed word, symbol, hypothesis or number of the counted statement or of its proof was altered. Layout and class: the article's own class is \texttt{article} with packages including neurips\_2024, inputenc, fontenc, hyperref, url, booktabs, amsfonts, nicefrac, microtype, xcolor, amsmath, amssymb, mathtools, amsthm; the wrapper is the lane's pinned \texttt{amsart} editorial wrapper at 10pt on A4, which restates the theorem-like environments the retained text needs and adds the article's own macro definitions that the retained text needs (none), with extra wrapper packages (bm, bbm, dsfont, xspace, cancel, mathtools). The delivered numbering is the unit's own. Assets: no retained span contains a figure environment, a graphics-inclusion command, a PDF-page inclusion, a table or a TikZ picture; no image file is copied into this package, the package declares no asset dependency and no image file is redistributed with it. Licence: version arXiv:2311.03087v3 of this article is distributed under the Creative Commons Attribution 4.0 International licence, as recorded in the arXiv licence metadata for this exact version and shown on the abstract page https://arxiv.org/abs/2311.03087v3. A full rights notice is delivered at licenses/arxiv-2311.03087-CC-BY-4.0.txt. Characters: every retained excerpt is pure ASCII, so the delivered engine, XeLaTeX with its default Latin Modern text font, reports no missing character.
\begin{lemma}\label{lem:pm}
    For any $m \in \mathbb{N}$, the map $p_m:\mathcal{D} \to \mathbb{R}_{\geq 0}$ assigning to a diagram its $m$-th largest persistence value is continuous.
\end{lemma}

\begin{proof}
    For a persistence diagram $D \in \mathcal D$ and a feature $f=(\tau_b, \tau_d) \in D$ we denote by $p(f)=\tau_d - \tau_b$ the persistence of $f$. Note that for two features $f,f'$ in a persistence diagram, we have $|p(f)- p(f')| \leq \sqrt{2}\|f -f'\|_\infty$.
    
    We need to show that for any $\varepsilon > 0$ and any $D \in \mathcal{D}$, we can choose some $\delta>0$ such that $|p_m(D') -p_m(D)| < \varepsilon$ for all $D'$ whose bottleneck distance from $D$ is at most $\delta$, i.e., $D'\in B_{\delta}(D)$. Recall that in this setting, by definition of the bottleneck distance, there is a bijection (of multisets) $\eta: D \to D'$ such that for all $f\in D$ we have $\|f-\eta(f)\|_{\infty} < \delta$. 

    Continuity of $p_m$ is easier to demonstrate at diagrams $D\in \mathcal{D}$ for which all off-diagonal points have different persistences. We discuss this case first. Assume $D$ has $l$ off-diagonal points. Denote by $f_k$ the feature with $k$-th largest persistence. 

    Let $\tilde{\delta} = \min\{p(f_k)-p(f_{k+1}) \mid k\leq l\}$ and $\delta = \min\big(\tilde{\delta} / (2\sqrt{2}), \varepsilon/(2\sqrt{2})\big)$. Consider $D'\in B_{\delta}(D)$.  By the choice of $\delta$, we have that $|p(\eta(f_k)) - p(f_k)| < \tilde{\delta}/2$, so that the values $p(\eta(f_1)),\dots, p(\eta(f_l)), 0$ are all distinct and strictly decreasing. As a result, for all $m\leq l$ we have $p_m(D') = p(\eta(f_m))$ and thus
    $$|p_m(D') - p_m(D)| = |p(\eta(f_m)) - p(f_m)| < \sqrt{2} \delta < \varepsilon.$$

    For $m>l$, we have $p_m(D)=0$ and any $f'\in D'$ with $p_m(D')=p(f')$ is paired with a point on the diagonal of $D$ under $\eta$. Hence
    $$|p_m(D') - p_m(D)| = |p(f') - 0|  < \sqrt{2} \delta < \varepsilon.$$

    The case where $D$ might contain off-diagonal points with identical persistence follows the same overall argument, but requires more careful bookkeeping. Assume $D$ has $l$ off-diagonal points. Let $\tilde{\delta} = \min \big(\{p_k(D) - p_{k+1}(D) \mid k\leq l\}\backslash \{0\}\big)$. Let $$\delta = \min\big(\tilde{\delta}/(2\sqrt{2}), \varepsilon /(2\sqrt{2})\big).$$
    Set $$F_k = \{f\in D \mid p(f)= p_k(D)\}$$ for $k \leq l$ and $F_{l+1}$ equal to the diagonal. These multisets all contain at least one element. We have $p_k(D) = p_{k+1}(D)$ if and only if $F_k = F_{k+1}$.

    As a result, for all $k\leq l$, we have $|\bigcup_{i=1}^{k} F_i |= k$ if and only if $F_k \neq F_{k+1}$ which in turn holds if and only if $p_k(D) \neq p_{k+1}(D)$.

    For any $f, \tilde{f} \in D$ with $p(f) > p(\tilde{f})$ we have 
    \begin{align}
        p(\eta(f))- p(\eta(\tilde{f})) &= p(\eta(f))-p(f) + p(f) - p(\tilde{f}) + p(\tilde{f}) - p(\eta(\tilde{f})) \nonumber \\
        &\geq p(f) - p(\tilde{f}) - | p(\eta(f))-p(f)| - |p(\tilde{f}) - p(\eta(\tilde{f}))| \nonumber \\
        &> p(f) - p(\tilde{f}) - \tilde{\delta}\nonumber \\
        &\geq 0,
    \end{align}
    so $p(\eta(f)) > p(\eta(\tilde{f}))$. We will now show that if $m\leq l$  and $f'\in D'$ such that $p_m(D') = p(f')$, then $\eta^{-1}(f') \in F_m$.

    Let $a$ be the highest number below $m$ such that $F_a\neq F_m$. Then 
    $$
    \left|\bigcup_{i=1}^a \eta(F_i)\right| = \left|\bigcup_{i=1}^a F_i\right| = a.
    $$
    and $\bigcup_{i=1}^a \eta(F_i)$ contains the $a<m$ features of $D'$ with smallest persistence. In particular, \mbox{$\eta^{-1}(f') \notin \bigcup_{i=1}^a F_i$.}

    Let $b$ be the largest number such that $F_h = F_m$. By a similar argument, $\bigcup_{i=1}^b \eta(F_i)$ contains the $b\geq m$ features in $D'$ with largest persistence, including $f'$. Thus,  $\eta^{-1}(f')\in \bigcup_{i=a+1}^{b} F_i$. But by choice of $a$ and $b$, we have $F_{a+1} = \dots = F_m=\dots =F_b$, so $\eta^{-1}(f')\in F_m$.

    Thus, 
    $$
    |p_m(D') - p_m(D)| = |p(f') - p(\eta^{-1}(f'))| \leq \sqrt{2} \|f'-\eta^{-1}(f')\|_{\infty} < \varepsilon.
    $$

    For $m> l$, we have $p_m(D)=0$ and $F_m$ is the set of points on the diagonal of $D$. Moreover, 
    $$\left|\bigcup_{i=1}^l F_i\right| = l = \left|\bigcup_{i=1}^l \eta(F_i)\right|,$$
    so that any $f'\in D'$ with $p_m(D') = p(f')$ must be in the image of the diagonal under $\eta$ and we can proceed as in the case of unique positive persistences.
\end{proof}

\end{document}
